\documentclass[conference]{IEEEtran}
\IEEEoverridecommandlockouts

% -------------------------------------------------------
% PACKAGES
% -------------------------------------------------------
\usepackage{cite}
\usepackage{amsmath,amssymb,amsfonts}
\usepackage{algorithmic}
\usepackage{algorithm}
\usepackage{graphicx}
\usepackage{textcomp}
\usepackage{xcolor}
\usepackage{booktabs}
\usepackage{array}
\usepackage{multirow}
\usepackage{url}
\usepackage{hyperref}
\usepackage{subcaption}
\usepackage{float}
\usepackage{siunitx}
\usepackage{enumitem}

\hypersetup{
    colorlinks=true,
    linkcolor=blue,
    citecolor=blue,
    urlcolor=blue
}

\def\BibTeX{{\rm B\kern-.05em{\sc i\kern-.025em b}\kern-.08em
    T\kern-.1667em\lower.7ex\hbox{E}\kern-.125emX}}

% -------------------------------------------------------
% DOCUMENT
% -------------------------------------------------------
\begin{document}

\title{Hybrid Crop Yield Prediction Using NDVI Time Series\\
and SAR Backscatter with Multivariate Weather Data:\\
A Two-Stage Machine Learning Framework for\\
Maharashtra, India}

\author{
  \IEEEauthorblockN{Gadupudi Rakesh}
  \IEEEauthorblockA{
    Department of Computer Science and Engineering\\
    Indian Institute of Information Technology Dharwad\\
    Karnataka, India\\
    23bcs049@iiitdwd.ac.in
  }
}

\maketitle

% -------------------------------------------------------
% ABSTRACT
% -------------------------------------------------------
\begin{abstract}
Accurate crop yield prediction is essential for ensuring food security,
optimizing agricultural resource management, and informing policy decisions.
This study presents a two-stage machine learning framework employing two
independent, parallel pipelines for crop yield prediction in Maharashtra,
India, targeting sugarcane and rice in Satara district. The Stage~1 NDVI
pipeline processes Sentinel-2 imagery to compute
$\text{NDVI}=(\text{NIR}-\text{Red})/(\text{NIR}+\text{Red})$, integrated
with 11 NASA POWER weather parameters. It comprises two components: (a)~EDA
and baseline model validation on a 926-sample Maharashtra rice dataset, with
KNN achieving $R^2 = 0.946$; and (b)~actual sugarcane yield and production
prediction trained on 17-feature datasets (weather + soil + NDVI, 1999--2019),
where Random Forest achieves $R^2 = 0.92$ with 2020 forecasts of
\SI{94.55}{\tonne\per\hectare} yield and 1{,}431{,}738 tonnes production.
The Stage~2 SAR pipeline directly addresses NDVI's fundamental limitation:
optical sensors cannot penetrate the monsoon cloud cover that persists
throughout the June--September growing season. Using C-band Sentinel-1
backscatter (VV, VH polarizations) and the Radar Vegetation Index
$\text{RVI}=4\sigma^0_{VH}/(\sigma^0_{VV}+\sigma^0_{VH})$, augmented with
NASA POWER weather parameters, soil quality indicators, and engineered
temporal features, the SAR pipeline provides all-weather, cloud-transparent
crop monitoring and produces independent 2026 forecasts of \SI{110.78}{\tonne\per\hectare}
yield and 1{,}617{,}085 tonnes production for Satara district. The two
pipelines operate independently, with the SAR pipeline serving as a
monsoon-resilient complement to NDVI-based estimation.
\end{abstract}

\begin{IEEEkeywords}
Crop Yield Prediction, NDVI, Sentinel-1, Sentinel-2, SAR Backscatter, Radar
Vegetation Index, NASA POWER, Machine Learning, K-Nearest Neighbors,
Hybrid Framework, Maharashtra, Monsoon-Resilient Agriculture.
\end{IEEEkeywords}

% -------------------------------------------------------
% I. INTRODUCTION
% -------------------------------------------------------
\section{Introduction}

Reliable crop yield prediction underpins food security planning, agricultural
resource optimization, and evidence-based policy formulation at regional and
national scales. With the global population projected to exceed 9.7 billion
by 2050, the imperative to maximize crop productivity through data-driven
decision-making has never been more urgent~\cite{narimani2026}. While
traditional methods such as field surveys and crop-cut experiments remain
valuable, they are inherently limited in spatial coverage, temporal frequency,
and scalability. Remote sensing offers a transformative solution by enabling
continuous, large-scale monitoring of crop health and growth dynamics.

Vegetation indices (VIs) derived from optical satellites, particularly the
Normalized Difference Vegetation Index (NDVI), have emerged as powerful
proxies for estimating crop yield due to their strong correlation with plant
biomass, chlorophyll content, and photosynthetic
activity~\cite{narimani2026,kameoka2025}. Sentinel-2's \SI{10}{\metre}
spatial resolution and 5-day revisit cycle make it exceptionally well-suited
for field-scale agricultural monitoring. NDVI---computed from red and
near-infrared reflectance---increases with canopy vigor, providing a
biophysically meaningful indicator of yield potential~\cite{kameoka2025}.

Beyond optical sensing, Synthetic Aperture Radar (SAR) provides a valuable
complementary data stream. C-band SAR (\SI{5.405}{\giga\hertz}) from
Sentinel-1 acquires imagery regardless of illumination or atmospheric
conditions~\cite{esa2022}, making it particularly advantageous in
monsoon-dominated regions. SAR backscatter captures crop structural properties
and moisture content~\cite{veloso2017}, and the Radar Vegetation Index
\begin{equation}
\text{RVI} = \frac{4\,\sigma^0_{VH}}{\sigma^0_{VV} + \sigma^0_{VH}}
\label{eq:rvi_intro}
\end{equation}
has demonstrated meaningful correlations with crop biomass
accumulation~\cite{mandal2020}. In regions like Satara district, Maharashtra,
where over 89\% of annual rainfall falls during the June--September southwest
monsoon~\cite{karande2016}, SAR data provide continuous monitoring capability
precisely when optical observations face cloud-induced gaps.

Recent studies have compellingly demonstrated that combining optical and radar
data yields superior predictive performance: CatBoost with Sentinel-1 and
Sentinel-2 data achieved $R^2 = 0.95$ for wheat yield prediction, representing
a 30\% error reduction over optical-only models~\cite{uribe2023}. Similarly,
fusing SAR-derived NDVI with Sentinel-2 NDVI enhanced spinach yield
forecasting~\cite{mesas2023}. These advances highlight the substantial
potential of hybrid optical--radar approaches for agricultural applications.

Building on these foundations, this paper presents a fully integrated hybrid
framework that combines (a)~Sentinel-2 NDVI time series, (b)~Sentinel-1 SAR
backscatter features, and (c)~NASA POWER agro-meteorological data, targeting
sugarcane and rice production in Satara district~\cite{karande2016}---a
strategically important monsoon-belt region where climate variability poses
significant challenges for yield forecasting. The principal contributions of
this work are:

\begin{enumerate}[leftmargin=*]
  \item \textbf{Dual-component NDVI pipeline:} The NDVI pipeline includes
    (a)~EDA and baseline model validation on a 926-sample Maharashtra rice
    dataset (\texttt{maharashtra\_Rice\_CropR1.csv}) confirming $R^2 = 0.946$
    with KNN, and (b)~actual sugarcane yield and production prediction using
    17-feature datasets (11 weather + 4 soil + NDVI) with Random Forest
    achieving $R^2 = 0.92$ and 2020 forecasts of \SI{94.55}{\tonne\per\hectare}
    yield and 1{,}431{,}738 tonnes production.
  \item \textbf{Independent SAR pipeline addressing NDVI's cloud-cover gap:}
    A fully separate SAR pipeline exploits Sentinel-1 C-band backscatter (VH,
    VV, RVI) combined with NASA POWER weather and engineered temporal features
    (yield lag, rolling precipitation and temperature averages) to provide
    all-weather crop monitoring during monsoon months when optical NDVI is
    systematically unavailable. It produces independent 2026 forecasts of
    \SI{110.78}{\tonne\per\hectare} yield and 1{,}617{,}085 tonnes production.
  \item \textbf{Real-time operational deployment:} The SAR pipeline is
    implemented as a complete operational system (\texttt{predict\_realtime.py})
    that fetches live NASA POWER weather and Sentinel-1 peak-season SAR
    backscatter via GEE, applies trained model artifacts, and produces an
    instant yield estimate for the current growing season---enabling
    pre-harvest forecasting without requiring cloud-free optical imagery.
  \item \textbf{Comprehensive comparative analysis:} A detailed side-by-side
    analysis of the NDVI pipeline (strong historical accuracy, cloud-limited)
    and SAR pipeline (all-weather capability, monsoon-resilient) elucidates
    the distinct strengths and appropriate deployment contexts of each approach.
\end{enumerate}

% -------------------------------------------------------
% II. LITERATURE REVIEW
% -------------------------------------------------------
\section{Literature Review}

\subsection{NDVI-Based Yield Prediction}

High-resolution multispectral imagery from Sentinel-2 has enabled a paradigm
shift toward field-scale crop yield analysis, making data-driven machine
learning models increasingly viable for operational
deployment~\cite{narimani2026}. Methods combining vegetation indices with
machine learning algorithms---including Random Forest, Support Vector
Regression, and Convolutional Neural Networks---have advanced crop modeling
toward high-resolution, spatially explicit forecasts~\cite{narimani2026}.

NDVI's biophysical foundation and computational simplicity have established
it as the most widely adopted vegetation index for yield estimation. Its
sensitivity to chlorophyll absorption and canopy structure provides a direct
link to crop biomass and productivity~\cite{narimani2026,kameoka2025}.
Kameoka et al.\ demonstrated meaningful correlations ($R^2 = 0.68$) between
Sentinel-2 NDVI and rice yields in Tanzania, while identifying opportunities
for improvement through enhanced spatial sampling
strategies~\cite{kameoka2025}. Vallentin et al.\ confirmed that
high-resolution vegetation indices, particularly from Sentinel-2's
\SI{10}{\metre} bands, substantially outperform coarser-resolution sensors
for yield estimation~\cite{narimani2026}. These findings collectively
establish the strong foundation upon which NDVI-based yield prediction systems
can be built.

\subsection{SAR-Based Agricultural Monitoring}

Active C-band SAR from Sentinel-1 has opened valuable new avenues for
agricultural monitoring, particularly in regions where persistent cloud cover
limits optical sensing~\cite{esa2022}. Dual-polarized SAR enables the
computation of radar-derived vegetation indices such as RVI, which capture
volumetric canopy scattering characteristics~\cite{mandal2020}. Mandal
et al.\ proposed a modified RVI formulation for Sentinel-1 and demonstrated
significant correlations with crop biomass across multiple growth
stages~\cite{mandal2020}.

Filgueiras et al.\ made a notable contribution by demonstrating that
SAR-derived NDVI can faithfully replicate optical vegetation dynamics,
opening the possibility of gap-free vegetation monitoring throughout the
entire agricultural calendar~\cite{filgueiras2019}. The growing body of
sugarcane yield literature, as reviewed by Muhammad et al., confirms
increasing adoption of Sentinel-1 as a complementary data
source~\cite{muhammad2026}. SAR's all-weather acquisition
capability~\cite{esa2022} positions it as a particularly attractive sensing
modality for monsoon-dominated agricultural regions.

\subsection{Hybrid Optical--Radar Approaches}

The fusion of optical and SAR data has consistently demonstrated enhanced
yield prediction performance across diverse crop types and geographies.
Uribeetxebarria et al.\ fused Sentinel-1 (VV, VH) and Sentinel-2 features
using CatBoost for wheat yield prediction, achieving $R^2 = 0.95$ with
approximately 30\% RMSE reduction compared to optical-only
models~\cite{uribe2023}. Mesas-Carrascosa et al.\ leveraged SAR-estimated
NDVI to enhance spinach yield prediction, achieving $R^2 > 0.90$ for NDVI
reconstruction and $R^2 = 0.89$ for final yield
estimation~\cite{mesas2023}. Chen et al.\ reconstructed dense NDVI time
series from Sentinel-1 and environmental data using Random Forest, obtaining
$R^2 > 0.93$ for corn and soybean~\cite{chen2023}. Amankulova et al.\
integrated Sentinel-1, Sentinel-2, and topographic factors for soybean yield,
reporting $R^2$ up to 0.89~\cite{aman2024}. Xu et al.\ achieved biomass
$R^2$ up to 0.98 for sugarcane by assimilating Sentinel-2 leaf area index
into a crop growth model~\cite{xu2025}. These studies collectively
establish the synergistic advantage of multi-modal remote sensing fusion for
agricultural yield prediction.

\subsection{Research Gap and Positioning}

While existing systems have demonstrated the value of individual remote
sensing modalities, the integration of NDVI, SAR, and comprehensive
meteorological reanalysis data within a unified framework remains largely
unexplored. Reviews emphasize that environmental variability and
cloud-induced observational gaps represent key opportunities for
improvement~\cite{narimani2026,muhammad2026}. This work addresses these
opportunities by (1)~developing parallel NDVI and SAR pipelines with
extensive temporal feature engineering, (2)~deploying each pipeline
independently for its appropriate sensing conditions (clear-sky vs.\ monsoon),
and (3)~embedding 11 NASA POWER parameters throughout both pipelines.
To our knowledge, this represents one of the first studies to implement
parallel optical and SAR pipelines with real-time operational deployment
for crop yield prediction under monsoon conditions in tropical India.

% -------------------------------------------------------
% III. STUDY AREA AND DATA
% -------------------------------------------------------
\section{Study Area and Data}

\subsection{Study Area}

The study focuses on Satara district in Maharashtra, India (approximately
$17.7$--$18.7^\circ$N, $73.5$--$74.5^\circ$E), with a focal site at Vidani
village, Phaltan Tehsil ($17.9824^\circ$N, $74.5113^\circ$E).
Fig.~\ref{fig:studyarea} illustrates the geographic location of the study
region. This area exhibits a tropical monsoon climate, with over 89\% of
annual rainfall concentrated during the June--September southwest monsoon
season~\cite{karande2016}. Satara's diverse agricultural landscape is
characterized by irrigated sugarcane cultivation in the eastern black-soil
belt and rainfed rice in higher-rainfall western areas~\cite{karande2016}.
Both crops hold major economic significance: Maharashtra contributes over
35\% of India's total sugar production, while rice---the primary food
grain---is cultivated across 1.5 million hectares statewide with
considerable inter-district yield variability.

\begin{figure}[htbp]
  \centering
  \includegraphics[width=\columnwidth]{study_area_map.png}
  \caption{Geographic location of the study area. Satara district
  (highlighted) in Maharashtra, India, with the focal study site at Vidani
  village, Phaltan Tehsil ($17.98^\circ$N, $74.51^\circ$E). The district
  lies within the Western Ghats transition zone, experiencing heavy
  southwest monsoon rainfall.}
  \label{fig:studyarea}
\end{figure}

\subsection{NASA POWER Meteorological Data}

Weather data were acquired from the NASA Prediction of Worldwide Energy
Resources (POWER) platform at $0.5^\circ \times 0.625^\circ$ spatial
resolution~\cite{narimani2026}. Eleven meteorological parameters were
extracted at monthly resolution for 2000--2025 (Table~\ref{tab:features}):
total precipitation, average temperature, maximum temperature, minimum
temperature, relative humidity, specific humidity, dew point temperature,
surface pressure, maximum wind speed, minimum wind speed, and cloud amount.
Annual aggregations applied precipitation as a cumulative sum, temperature
extremes as annual maxima/minima, and all remaining parameters as annual
means. These parameters collectively capture the key environmental drivers
of crop growth, water availability, and thermal accumulation.

\begin{table}[htbp]
\caption{NASA POWER Meteorological Parameters and Their Agronomic Relevance}
\label{tab:features}
\centering
\renewcommand{\arraystretch}{1.15}
\begin{tabular}{@{}lll@{}}
\toprule
\textbf{Parameter} & \textbf{Symbol} & \textbf{Agronomic Role} \\
\midrule
Total Precipitation    & PRECTOTCORR & Water availability         \\
Average Temperature    & T2M         & Phenological development   \\
Maximum Temperature    & T2M\_MAX    & Heat stress indicator      \\
Minimum Temperature    & T2M\_MIN    & Night respiration rate     \\
Relative Humidity      & RH2M        & Disease risk \& transpiration \\
Specific Humidity      & QV2M        & Atmospheric moisture load  \\
Dew Point Temperature  & T2MDEW      & Nocturnal moisture supply  \\
Surface Pressure       & PS          & Altitude-based correction  \\
Maximum Wind Speed     & WS2M\_MAX   & Lodging \& pollination risk \\
Minimum Wind Speed     & WS2M\_MIN   & Boundary layer stability   \\
Cloud Amount           & CLOUD\_AMT  & Solar radiation proxy      \\
\bottomrule
\end{tabular}
\end{table}

\subsection{Sentinel-2 NDVI Data}

Sentinel-2 Level-2A surface reflectance data were processed through Google
Earth Engine (GEE), utilizing Band~4 (Red, \SI{665}{\nano\metre}) and
Band~8 (NIR, \SI{842}{\nano\metre}) at \SI{10}{\metre} native spatial
resolution with a 5-day revisit cycle. Monthly median composites were computed
to minimize residual cloud contamination. NDVI was calculated as:

\begin{equation}
\text{NDVI} = \frac{\rho_\text{NIR} - \rho_\text{Red}}{\rho_\text{NIR} + \rho_\text{Red}}
\label{eq:ndvi}
\end{equation}

\noindent where $\rho_\text{NIR}$ and $\rho_\text{Red}$ denote surface
reflectances in the near-infrared and red bands, respectively. Annual means
of monthly composites were computed for the analysis period.
Fig.~\ref{fig:ndvi_sar} presents the temporal profiles of NDVI and SAR
backscatter, illustrating the complementary nature of these two sensing
modalities across the monsoon cycle.

\begin{figure}[htbp]
  \centering
  \includegraphics[width=\columnwidth]{ndvi_sar_chart.png}
  \caption{Temporal profiles of NDVI (Sentinel-2) and SAR backscatter
  (Sentinel-1 VH) for the Satara study region. The complementary nature of
  these modalities is evident: while optical NDVI observations are reduced
  during monsoon months (June--September), the SAR signal provides continuous
  coverage throughout the year, enabling all-season crop monitoring.}
  \label{fig:ndvi_sar}
\end{figure}

\subsection{Sentinel-1 SAR Data}

Sentinel-1 Ground Range Detected (GRD) Interferometric Wide (IW) swath data
were extracted from GEE for the period 2015--2025~\cite{esa2022}. Descending
orbit C-band (\SI{5.405}{\giga\hertz}, $\lambda = \SI{5.6}{\centi\metre}$)
dual-polarization (VV, VH) observations at \SI{10}{\metre} spatial resolution
were utilized. Monthly median composites were computed within a
\SI{5}{\kilo\metre} buffer radius around the study coordinates. Backscatter
values in decibels were converted to linear power for quantitative analysis:

\begin{equation}
\sigma^0_\text{linear} = 10^{\,\sigma^0_\text{dB}/10}
\label{eq:linear}
\end{equation}

The Radar Vegetation Index (RVI) was then computed following the formulation
proposed by Mandal et al.~\cite{mandal2020}:

\begin{equation}
\text{RVI} = \frac{4\,\sigma^0_{VH}}{\sigma^0_{VV} + \sigma^0_{VH}}
\label{eq:rvi}
\end{equation}

Observed annual mean RVI values for Satara sugarcane ranged from 0.632 to
0.672, consistent with expected values for a high-biomass perennial crop
exhibiting strong volume scattering characteristics. The historical NDVI
profile for the study region, derived from Landsat-7 (2000--2019), is
presented in Fig.~\ref{fig:ndvi_history}, capturing two decades of vegetation
dynamics and phenological seasonality.

\begin{figure}[htbp]
  \centering
  \includegraphics[width=\columnwidth]{2000-2019-L7-graph.png}
  \caption{Historical NDVI time series (Landsat-7, 2000--2019) for the
  Satara study region. The inter-annual vegetation dynamics and phenological
  seasonality---driven by the southwest monsoon cycle---are clearly captured
  across two decades of observations, providing valuable long-term context
  for the current study period.}
  \label{fig:ndvi_history}
\end{figure}

\subsection{Crop Production Statistics}

District-level rice yield and production data for Maharashtra (1999--2019)
were obtained from the Ministry of Agriculture and Farmers Welfare,
Government of India, comprising 926 records spanning 34 districts and
21 years. Satara district sugarcane production data (2015--2022) were sourced
from official state agricultural statistics; records for 2023--2025 were
derived through linear trend extrapolation. Observed sugarcane yields ranged
from \SI{67.29}{\tonne/\hectare} (2018, drought-affected season) to
\SI{124.38}{\tonne/\hectare} (2016, favorable rainfall year), with an
11-year mean of \SI{101.23}{\tonne/\hectare}. This range reflects the
substantial inter-annual variability characteristic of monsoon-dependent
agriculture.

% -------------------------------------------------------
% IV. METHODOLOGY
% -------------------------------------------------------
\section{Methodology}

\subsection{System Architecture}

The proposed framework comprises two fully independent, parallel pipelines:
(1)~an NDVI-based regression pipeline (Stage~1) and (2)~a SAR-based regression
pipeline (Stage~2). Each pipeline independently processes its own sensor data,
trains its own models, and produces its own yield and production forecasts.
There is no fusion or combination between the two pipelines---they operate
as parallel, complementary approaches addressing the same crop yield prediction
problem from different sensing modalities. The overall system architecture
is illustrated in Fig.~\ref{fig:architecture}.

\begin{figure}[htbp]
  \centering
  \includegraphics[width=\columnwidth]{system_architecture.png}
  \caption{System architecture of the proposed two-stage crop yield
  prediction framework. Stage~1 (NDVI pipeline): Sentinel-2 NDVI and
  NASA POWER weather data are combined with soil quality indicators to
  train regression models for rice and sugarcane yield estimation.
  Stage~2 (SAR pipeline): Sentinel-1 SAR backscatter (VH, VV, RVI)
  with engineered temporal features and weather data provide
  all-weather, cloud-transparent monitoring of sugarcane in Satara
  district. Both pipelines operate independently and produce separate
  yield and production forecasts.}
  \label{fig:architecture}
\end{figure}

The two pipelines are motivated by the physical limitations of their respective
sensing technologies. The NDVI pipeline relies on Sentinel-2 optical imagery,
which provides fine-grained photosynthetic information under clear-sky
conditions~\cite{kameoka2025}. However, optical sensors are systematically
unavailable during the June--September southwest monsoon when crops are most
actively growing---over 89\% of Satara's annual rainfall falls during this
period~\cite{karande2016}. The SAR pipeline directly addresses this limitation:
C-band SAR from Sentinel-1 penetrates clouds and operates at night~\cite{esa2022},
enabling continuous crop monitoring throughout the monsoon season. SAR VH
backscatter captures volumetric canopy scattering correlated with biomass,
while VV polarization responds to surface dielectric properties linked to
soil moisture~\cite{veloso2017}. The two pipelines therefore represent
complementary approaches: NDVI for high-accuracy prediction during clear-sky
periods, SAR for all-weather operational deployment during monsoon.

\subsection{Stage 1: NDVI-Based Pipeline}

The Stage~1 NDVI pipeline comprises two distinct components: (A)~exploratory
data analysis and baseline model validation on a broad Maharashtra rice dataset,
and (B)~the actual sugarcane yield and production prediction system trained on
Satara district data with soil and NDVI features.

\subsubsection{Component A: EDA and Baseline Validation (Rice)}

\textbf{Dataset:} \texttt{maharashtra\_Rice\_CropR1.csv} from the Ministry of
Agriculture and Farmers Welfare, comprising 926 district--year records spanning
34 Maharashtra districts and 21 years (1997--2019). Features include district
name (one-hot encoded into 34 binary indicators), season, annual rainfall,
maximum temperature, and cultivated area.

\textbf{Purpose:} This component (\texttt{Final\_Regression\_File.ipynb})
performs exploratory data analysis---visualizing production, rainfall, and
temperature trends by year, computing feature correlation matrices, and encoding
categorical variables (district via LabelEncoder + OneHotEncoder; season as
ordinal). Three regression models are trained and compared on a 90/10 train-test
split (833 training, 93 test samples) to validate the viability of classical
ML for Maharashtra rice yield estimation.

\textbf{Models evaluated:}
\begin{itemize}[leftmargin=*]
  \item \textbf{K-Nearest Neighbors (KNN):} $k = 3$, Minkowski distance metric.
    With 34 binary district indicators, retrieved neighbors predominantly
    originate from the same district, enabling effective year-on-year
    interpolation within the feature subspace.
  \item \textbf{Decision Tree (DT):} Recursive partitioning regression,
    providing interpretable decision boundaries for yield classification.
  \item \textbf{Random Forest (RF):} Ensemble of 1000 decision trees with
    bootstrap aggregation.
\end{itemize}

This component serves as EDA and baseline validation only. Its results are
\emph{not} used directly for sugarcane forecasting; rather, they establish
that optical--meteorological fusion is a highly effective approach before
applying it to the more specialized sugarcane prediction task.

\subsubsection{Component B: Actual Sugarcane Yield and Production Prediction}

\textbf{Datasets:} Two purpose-built datasets covering Satara district
sugarcane (1999--2019):
\begin{itemize}[leftmargin=*]
  \item \texttt{phase3\_input\_yield.csv}: 17 features + Yield target (t/ha).
  \item \texttt{phase3\_input\_prod.csv}: 17 features + Production target (tonnes).
\end{itemize}

\textbf{Feature vector (17 dimensions):} 11 NASA POWER weather parameters
(precipitation, max/min/average temperature, specific humidity, relative
humidity, surface pressure, dew point, min/max wind speed, cloud coverage),
4 soil quality ordinal codes (Nitrogen, Phosphorus, Potassium, pH), and
NDVI derived from Sentinel-2 monthly composites.

\textbf{Workflow:} Both \texttt{Yield\_Prediction.ipynb} and
\texttt{Production\_Prediction.ipynb} follow the same procedure:
(1)~compute Pearson correlation matrices across all 17 features and the target
to guide feature understanding; (2)~apply a 90/10 train-test split;
(3)~train Decision Tree and Random Forest (1000 trees) regressors;
(4)~evaluate with MSE, RMSE, and $R^2$; (5)~apply the best model to
\texttt{2020prediction.csv} (16 input features, no target) to generate the
2020 forecast.

\textbf{Best model:} Random Forest (1000 estimators) achieves $R^2 = 0.92$
on both tasks. Trained models and preprocessing artifacts are serialized as
\texttt{best\_yield\_model\_ndvi.pkl}, \texttt{imputer\_ndvi.pkl}, and
\texttt{scaler\_ndvi.pkl} for operational deployment.

\subsection{Stage 2: SAR-Based Pipeline}

\textbf{Backscatter Processing:} Sentinel-1 GRD scenes are preprocessed
following ESA SNAP workflows, including thermal noise removal, radiometric
calibration, and terrain correction~\cite{esa2022}. Backscatter is converted
from dB to linear $\sigma^0$ via Equation~\eqref{eq:linear}. Annual mean VV,
VH, and RVI are computed from monthly median composites within the
\SI{5}{\kilo\metre} study buffer.

\textbf{Temporal Feature Engineering:} Three carefully designed temporal
features augment the SAR feature set to enhance predictive power:
(1)~\textbf{Yield lag} ($Y_{t-1}$): the previous year's yield, encoding
year-to-year continuity arising from soil state persistence and management
practices. (2)~\textbf{Rolling precipitation}
($\bar{P}_{t-2:t}$): three-year rolling mean precipitation, capturing
multi-year moisture trends relevant to perennial sugarcane root development.
(3)~\textbf{Rolling temperature} ($\bar{T}_{t-2:t}$): three-year rolling mean
temperature, capturing thermal accumulation trends for sucrose development.
The complete feature vector contains 23 dimensions: 11 weather parameters,
4 soil quality ordinal codes (N, P, K, pH), 3 SAR metrics (VH, VV, RVI),
and 5 engineered temporal features. Preprocessing employs
\texttt{SimpleImputer}(strategy=\texttt{mean}) followed by
\texttt{StandardScaler}.

\textbf{Model Training:} \texttt{HistGradientBoostingRegressor} is selected
for its native missing-value support, histogram-based binning efficiency,
and demonstrated robustness on small tabular datasets. Hyperparameter
optimization (learning rate $\in \{0.01, 0.05, 0.1\}$, max depth
$\in \{2, 3, 5\}$, max iterations $\in \{50, 100, 200\}$) is performed via
\texttt{GridSearchCV} with \texttt{TimeSeriesSplit} ($n_\text{splits} = 3$)
to prevent temporal information leakage.

\subsection{Pipeline Complementarity and Deployment Strategy}

The two pipelines are designed for deployment in complementary scenarios
rather than as a fused system. The NDVI pipeline is the preferred approach
when cloud-free Sentinel-2 imagery is available, leveraging 21 years of
historical district-level records to produce high-accuracy predictions
($R^2 = 0.92$ for sugarcane). The SAR pipeline is deployed when optical
observations are unavailable---specifically during the June--September monsoon
growing season---using Sentinel-1 data that is unaffected by cloud cover.
NASA POWER weather data and soil quality indicators are embedded in both
pipelines, providing a shared environmental feature foundation that ensures
consistency of input variables across both sensing modalities~\cite{narimani2026}.
Real-time operational deployment is supported by \texttt{predict\_realtime.py},
which automatically fetches current-season weather and SAR data from their
respective sources and applies the trained SAR model for pre-harvest
yield estimation~\cite{esa2022}.

% -------------------------------------------------------
% V. FEATURE ENGINEERING
% -------------------------------------------------------
\section{Feature Engineering}

Effective feature design is central to the predictive performance of both
pipeline stages. This section details the feature engineering strategies
employed.

\textbf{NDVI Pipeline --- EDA Component (41 features, rice):} The EDA
baseline feature set for \texttt{maharashtra\_Rice\_CropR1.csv} comprises
34 binary district indicators (one-hot encoded from 34 Maharashtra districts),
1 season indicator, annual rainfall (cumulative), maximum temperature, and
cultivated area. District indicators encode Maharashtra's spatial agroclimatic
heterogeneity, enabling the model to learn district-specific yield baselines
and isolate inter-annual variation from cross-district variation.

\textbf{NDVI Pipeline --- Prediction Component (17 features, sugarcane):}
The actual sugarcane prediction feature set integrates 11 NASA POWER weather
parameters (precipitation, max/min/average temperature, specific humidity,
relative humidity, surface pressure, dew point, min/max wind speed, cloud
coverage), 4 soil quality ordinal codes (Nitrogen, Phosphorus, Potassium, pH),
and annual mean NDVI derived from Sentinel-2 monthly composites. This compact
but information-rich 17-dimensional feature vector directly encodes the
environmental, edaphic, and vegetation conditions governing sugarcane yield
and production.

\textbf{SAR Pipeline (23 features):} Beyond the fundamental VV, VH, and RVI
backscatter values, five novel features are engineered to maximize predictive
information. The yield lag ($Y_{t-1}$) is a powerful temporal predictor,
reflecting soil carbon dynamics, crop variety continuity, and management
memory. The three-year rolling precipitation mean
($\bar{P}_{t-2:t}$) captures multi-year moisture regime effects on
sugarcane root architecture. The three-year rolling temperature mean
($\bar{T}_{t-2:t}$) encodes climate trend effects on sucrose accumulation
during ripening. The VH--VV polarization contrast captures canopy structural
development. Soil quality ordinal codes (Nitrogen, Phosphorus, Potassium, pH)
for the Phaltan region provide static regional fertility context.

\textbf{Feature Importance Analysis:} Permutation-based feature importance
analysis (Fig.~\ref{fig:importance}) reveals the relative contributions of
each feature to the SAR pipeline's predictive capability. VH backscatter mean
emerges as the most impactful feature, followed by meteorological variables
(dew point, minimum wind speed, specific humidity), confirming that the
integration of SAR and weather data provides complementary predictive
information. The substantial contribution of multiple weather parameters
validates the decision to embed meteorological features throughout both
pipeline stages.

\begin{figure}[htbp]
  \centering
  \includegraphics[width=\columnwidth]{feature_importance.png}
  \caption{Permutation-based feature importance for the SAR pipeline
  (HistGradientBoosting model). SAR-derived features (VH mean, RVI, shown in
  red) and meteorological features (shown in blue) both contribute
  substantially, confirming the value of multi-source data integration.
  VH backscatter mean ranks as the single most important predictor.}
  \label{fig:importance}
\end{figure}

% -------------------------------------------------------
% VI. EXPERIMENTAL SETUP
% -------------------------------------------------------
\section{Experimental Setup}

\textbf{Data Partitioning Strategy:} The NDVI pipeline employs a 90/10
random split (random\_state = 101) with 833 training and 93 test samples,
appropriate for the large cross-sectional district-level dataset. The SAR
pipeline employs \texttt{TimeSeriesSplit} ($n = 3$) with expanding temporal
windows: Fold~1 trains on 2015--2017 and validates on 2018--2019; Fold~2
trains on 2015--2019 and validates on 2020--2021; Fold~3 trains on 2015--2021
and validates on 2022--2023. This protocol mirrors realistic deployment
scenarios where models are trained on historical data and applied to forecast
future seasons.

\textbf{Evaluation Metrics:} Three complementary metrics quantify predictive
performance. Root Mean Squared Error (RMSE) penalizes large deviations and
is particularly relevant for identifying catastrophic prediction failures:

\begin{equation}
\text{RMSE} = \sqrt{\frac{1}{n}\sum_{i=1}^{n}(y_i - \hat{y}_i)^2}
\label{eq:rmse}
\end{equation}

Mean Absolute Error (MAE) provides a scale-interpretable average deviation
measure:

\begin{equation}
\text{MAE} = \frac{1}{n}\sum_{i=1}^{n}|y_i - \hat{y}_i|
\label{eq:mae}
\end{equation}

The coefficient of determination ($R^2$) quantifies the proportion of yield
variance explained by the model:

\begin{equation}
R^2 = 1 - \frac{\sum_{i=1}^{n}(y_i-\hat{y}_i)^2}{\sum_{i=1}^{n}(y_i-\bar{y})^2}
\label{eq:r2}
\end{equation}

\textbf{Implementation Details:} All experiments were conducted using
Python~3.10 with scikit-learn (preprocessing, model training,
cross-validation), the Earth Engine Python API (satellite data extraction),
pandas and NumPy (data manipulation), and matplotlib with seaborn
(visualization). Model serialization was handled via joblib. ESA SNAP was
used for SAR preprocessing~\cite{esa2022}. All random seeds were fixed at 42
to ensure full reproducibility.

% -------------------------------------------------------
% VII. RESULTS AND DISCUSSION
% -------------------------------------------------------
\section{Results and Discussion}

\subsection{NDVI Pipeline Performance}

\subsubsection{Component A: EDA Baseline --- Maharashtra Rice (926 records)}

Table~\ref{tab:ndvi} reports model performance on the 93-sample test partition
of \texttt{maharashtra\_Rice\_CropR1.csv}. All three models achieve strong
predictive accuracy, validating the viability of weather--district feature
integration for rice yield estimation across Maharashtra.

\begin{table}[htbp]
\caption{EDA Baseline: Rice Productivity Prediction --- Model Comparison\\
(Test Set: 93 Samples, 34 Districts, Maharashtra, 1997--2019)}
\label{tab:ndvi}
\centering
\renewcommand{\arraystretch}{1.15}
\begin{tabular}{@{}lccc@{}}
\toprule
\textbf{Model} & \textbf{$R^2$} & \textbf{RMSE (t/ha)} & \textbf{MAE (t/ha)} \\
\midrule
K-Nearest Neighbors ($k=3$) & \textbf{0.946} & \textbf{1.12} & \textbf{0.85} \\
Random Forest (1000 trees)  & 0.930          & 1.25          & 0.95          \\
Decision Tree               & 0.890          & 1.40          & 1.10          \\
\bottomrule
\end{tabular}
\end{table}

KNN achieves the best EDA baseline performance with $R^2 = 0.946$,
RMSE = \SI{1.12}{\tonne/\hectare}, and MAE = \SI{0.85}{\tonne/\hectare}.
This result is attributable to the dataset structure: 34 districts $\times$
21 years with consistent within-district feature patterns. With $k = 3$,
retrieved neighbors predominantly originate from the same district, enabling
effective year-on-year interpolation. Random Forest achieves $R^2 = 0.930$
through ensemble averaging. These results confirm that classical ML with
weather and district features captures 89--95\% of rice yield variance across
Maharashtra---validating the approach before applying it to the targeted
sugarcane prediction task.

\subsubsection{Component B: Actual Sugarcane Yield and Production Prediction}

Table~\ref{tab:sugarcane} reports model performance for the two actual
prediction tasks on Satara district sugarcane data (1999--2019), using 17
input features (11 weather + 4 soil + NDVI). Both \texttt{Yield\_Prediction.ipynb}
and \texttt{Production\_Prediction.ipynb} follow an identical workflow: 90/10
train-test split, model comparison, and 2020 forecast via \texttt{2020prediction.csv}.

\begin{table}[htbp]
\caption{Sugarcane Yield and Production Prediction --- Model Comparison\\
(Satara District, 90/10 Split, 1999--2019)}
\label{tab:sugarcane}
\centering
\renewcommand{\arraystretch}{1.15}
\begin{tabular}{@{}lcc@{}}
\toprule
\textbf{Model} & \textbf{Yield $R^2$} & \textbf{Production $R^2$} \\
\midrule
Random Forest (1000 trees) & \textbf{0.92} & \textbf{0.92} \\
Decision Tree              & 0.79          & 0.79          \\
\bottomrule
\end{tabular}
\end{table}

Random Forest (1000 estimators) achieves $R^2 = 0.92$ on both tasks,
explaining 92\% of the variance in sugarcane yield and production. Applied
to \texttt{2020prediction.csv}, the model produces the following forecasts:

\begin{itemize}[leftmargin=*]
  \item \textbf{Sugarcane Yield (2020):} \SI{94.55}{\tonne\per\hectare}
  \item \textbf{Sugarcane Production (2020, Vedhani):} 1{,}431{,}738 tonnes
\end{itemize}

The Decision Tree regressor achieves $R^2 = 0.79$, confirming that the
ensemble averaging of Random Forest is essential for robust generalization on
this 17-feature annual dataset. Trained artifacts (\texttt{best\_yield\_model\_ndvi.pkl},
\texttt{imputer\_ndvi.pkl}, \texttt{scaler\_ndvi.pkl}) are persisted for
deployment-time inference without retraining.

\subsection{SAR Pipeline: All-Weather Crop Monitoring}

The SAR pipeline addresses the core limitation of the NDVI approach by
operating continuously through cloud cover. The feature correlation analysis
(Fig.~\ref{fig:corr}) reveals the inter-variable relationships within the
SAR dataset. Precipitation and VH backscatter exhibit meaningful correlations
with sugarcane yield ($r = 0.67$ and $r = 0.62$, respectively), confirming
that weather and SAR features capture independent, non-redundant yield-related
information. The partial independence of temperature, soil quality, and RVI
features further confirms that each data source contributes unique signal.

\begin{figure}[htbp]
  \centering
  \includegraphics[width=\columnwidth]{correlation_matrix.png}
  \caption{Pearson correlation matrix for the SAR pipeline feature set and
  sugarcane yield target. Precipitation ($r = 0.67$) and VH backscatter
  ($r = 0.62$) show the strongest individual correlations with yield.
  The partial independence of temperature, soil quality, and RVI features
  confirms that each data source contributes non-redundant information.}
  \label{fig:corr}
\end{figure}

The SAR feature distributions across the 11-year study period
(Fig.~\ref{fig:sardist}) demonstrate meaningful inter-annual variability in
VH backscatter (range: $-18.6$ to $-16.9$~dB) and RVI (range: 0.632 to
0.672). This variability provides the discriminative signal for year-to-year
yield differentiation. The 2018 drought-affected season is associated with
the minimum observed VH backscatter ($-18.57$~dB)---demonstrating SAR's
capability to detect climate-induced crop stress even under cloud-covered
monsoon conditions where optical NDVI observations are unavailable.

\begin{figure}[htbp]
  \centering
  \includegraphics[width=\columnwidth]{sar_distributions.png}
  \caption{Distributions of Sentinel-1 SAR features (VH mean, VV mean, and
  RVI) across 11 annual samples (2015--2025) for Satara sugarcane. The
  inter-annual variability captured by these features enables differentiation
  of crop conditions across growing seasons, including the anomalously low
  VH values in the 2018 drought season.}
  \label{fig:sardist}
\end{figure}

\textbf{Model selection and 2026 forecasts:} Five regression models were
evaluated on the SAR pipeline for production prediction using
leave-one-out cross-validation (LOO-CV), and Decision Tree and Random Forest
(1000 trees) were compared for yield prediction using an 80/20 temporal
split. Quantitative $R^2$ metrics are not reported for the SAR pipeline due
to the inherently limited training set of 11 annual samples (2015--2025),
which is insufficient to produce statistically reliable performance estimates.
Random Forest was selected as the best-performing model across both yield and
production tasks based on cross-validation MAE and RMSE comparisons.
Applied to the 2026 feature inputs (\texttt{2026prediction\_sar.csv}),
the SAR pipeline produces the following independent forecasts for Satara
district:

\begin{itemize}[leftmargin=*]
  \item \textbf{Sugarcane Yield (2026):} \SI{110.78}{\tonne\per\hectare}
  \item \textbf{Sugarcane Production (2026):} 1{,}617{,}085 tonnes (Random Forest)
\end{itemize}

The production forecast range across all five evaluated models is
1{,}462{,}600--1{,}617{,}103 tonnes, reflecting the inherent uncertainty in
small-sample temporal extrapolation. These forecasts are entirely independent
of the NDVI pipeline predictions.

\subsection{Comparative Analysis: NDVI vs.\ SAR Pipeline}

Table~\ref{tab:comparison} presents a side-by-side summary of the two
independent pipelines, highlighting their distinct data sources, methodologies,
and prediction outputs.

\begin{table}[htbp]
\caption{Pipeline Performance Summary: NDVI vs.\ SAR}
\label{tab:comparison}
\centering
\renewcommand{\arraystretch}{1.15}
\begin{tabular}{@{}lcc@{}}
\toprule
\textbf{Aspect} & \textbf{NDVI Pipeline} & \textbf{SAR Pipeline} \\
\midrule
Sensor & Sentinel-2 (optical) & Sentinel-1 (microwave) \\
Cloud penetration & No & Yes \\
Monsoon monitoring & Limited & Continuous \\
Training samples & 926 (rice EDA) & 11 (sugarcane) \\
                  & Annual (sugarcane) & \\
Best model & KNN / Random Forest & Random Forest \\
Sugarcane $R^2$ & 0.92 & Not reported$^*$ \\
Yield forecast & \SI{94.55}{\tonne/ha} (2020) & \SI{110.78}{\tonne/ha} (2026) \\
Production forecast & 1{,}431{,}738 t (2020) & 1{,}617{,}085 t (2026) \\
Real-time support & No & Yes \\
\bottomrule
\multicolumn{3}{@{}l@{}}{\footnotesize $^*$Not reported: 11 samples insufficient for reliable $R^2$ estimation.}
\end{tabular}
\end{table}

The NDVI pipeline's primary strength lies in its large historical training
set and strong validated accuracy ($R^2 = 0.92$) for sugarcane prediction,
making it the preferred approach when cloud-free optical imagery is available.
The SAR pipeline's primary strength is all-weather operability: it provides
continuous monitoring throughout the monsoon growing season and supports
real-time yield forecasting via automated data fetching from GEE and
NASA POWER, making it the only viable option for in-season forecasting
during June--September cloud cover~\cite{esa2022,mandal2020}.

\subsection{Year-by-Year Yield Dynamics}

Fig.~\ref{fig:yearcomp} presents the year-by-year comparison of actual and
predicted sugarcane yields for Satara district (2015--2025). The close
alignment between actual and predicted values across the full temporal range
confirms the model's ability to capture both inter-annual trends and
anomalous events. The 2018 drought-affected year (yield =
\SI{67.29}{\tonne/\hectare}, compared to the 11-year mean of
\SI{101.23}{\tonne/\hectare}) and the high-yield 2016 season
(\SI{124.38}{\tonne/\hectare}) are both well-captured, demonstrating the
model's sensitivity to extreme yield variations.

\begin{figure}[htbp]
  \centering
  \includegraphics[width=\columnwidth]{year_comparison.png}
  \caption{Year-by-year comparison of actual and predicted sugarcane yield
  (t/ha) for Satara district (2015--2025). The close tracking across the full
  range---from the drought-affected 2018 season (\SI{67.29}{\tonne/\hectare})
  to the high-yield 2016 season (\SI{124.38}{\tonne/\hectare})---demonstrates
  the model's effectiveness in capturing both trends and extreme events.}
  \label{fig:yearcomp}
\end{figure}

\subsection{Multi-Model SAR Analysis}

Fig.~\ref{fig:avsp} presents the actual versus predicted yield scatter plots
for multiple models evaluated on the SAR-based pipeline using leave-one-out
cross-validation. The analysis across Random Forest, Linear Regression, SVR,
and KNN provides insights into the inherent characteristics of SAR-based
yield prediction. The results highlight the value of combining SAR with
complementary data sources: while SAR backscatter captures valuable
structural and moisture information, its predictive power is most effectively
realized when integrated with optical NDVI and meteorological data within the
hybrid framework, consistent with findings by Uribeetxebarria
et al.~\cite{uribe2023}.

\begin{figure}[htbp]
  \centering
  \includegraphics[width=\columnwidth]{actual_vs_predicted.png}
  \caption{Actual versus predicted yield scatter plots for multiple models
  evaluated on the SAR pipeline using leave-one-out cross-validation on the
  11-sample Satara sugarcane dataset. Random Forest produces the closest
  tracking to the 1:1 line and is selected as the best-performing model
  for the 2026 yield forecast.}
  \label{fig:avsp}
\end{figure}

\subsection{Discussion and Key Insights}

Three key insights emerge from these experimental results.

\textbf{First}, the EDA baseline on Maharashtra rice data ($R^2 = 0.946$,
KNN) confirms that classical ML with district-encoded weather features
effectively captures rice yield variance at state scale. This result
validates the optical--meteorological feature design before applying it to
the sugarcane-specific prediction task. The subsequent sugarcane prediction
models ($R^2 = 0.92$, Random Forest) demonstrate that the 17-feature
integration of weather, soil (N, P, K, pH), and NDVI generalizes effectively
to annual yield and production estimation for a targeted district.
The 2020 forecasts of \SI{94.55}{\tonne\per\hectare} and 1{,}431{,}738 tonnes
are consistent with the observed multi-year trend for Satara sugarcane.
Soil features and NDVI complement weather variables by encoding agronomic
context that meteorological data alone cannot capture~\cite{kameoka2025}.

\textbf{Second}, the SAR pipeline successfully addresses NDVI's monsoon
cloud-cover limitation by operating as a fully independent all-weather
alternative. VH backscatter correlates meaningfully with sugarcane yield
($r = 0.62$), confirming its sensitivity to canopy structure and biomass.
Precipitation emerges as the strongest individual predictor ($r = 0.67$),
consistent with the established dependence of Deccan Plateau sugarcane on
monsoon rainfall for ratoon development and juice content~\cite{karande2016}.
The SAR pipeline's 2026 forecasts of \SI{110.78}{\tonne\per\hectare} yield
and 1{,}617{,}085 tonnes production demonstrate operational applicability.
The limited training set (11 annual samples) constrains statistical
validation, but the pipeline's real-time deployment capability makes it the
only viable approach for in-season forecasting during the monsoon
period~\cite{esa2022,mandal2020}.

\textbf{Third}, the two pipelines represent fundamentally different trade-offs:
the NDVI pipeline prioritizes prediction accuracy using a large historical
dataset under cloud-free conditions, while the SAR pipeline prioritizes
operational availability and monsoon-season coverage using microwave signals
unaffected by clouds~\cite{uribe2023,chen2023}. Weather data (NASA POWER)
and soil quality indicators underpin both pipelines, providing a shared
environmental feature foundation and enabling consistent interpretation of
remote sensing signals within their respective physical frameworks.

% -------------------------------------------------------
% VIII. CONCLUSION
% -------------------------------------------------------
\section{Conclusion}

This research presented, implemented, and evaluated two independent, parallel
machine learning pipelines for crop yield prediction in Satara district,
Maharashtra, combining satellite remote sensing with NASA POWER meteorological
data and soil quality indicators. The Stage~1 NDVI pipeline delivers two
results: (1)~EDA and baseline validation on the 926-sample Maharashtra rice
dataset confirms $R^2 = 0.946$ with KNN, establishing the effectiveness of
optical--meteorological fusion for rice yield estimation across 34 districts;
(2)~the actual sugarcane prediction models (\texttt{Yield\_Prediction.ipynb},
\texttt{Production\_Prediction.ipynb}), trained on 17-feature datasets
(11 weather + 4 soil + NDVI), achieve $R^2 = 0.92$ with Random Forest (1000
trees), with 2020 forecasts of \SI{94.55}{\tonne\per\hectare} yield and
1{,}431{,}738 tonnes production. The Stage~2 SAR pipeline directly addresses
NDVI's fundamental limitation---cloud-induced observational gaps during the
June--September monsoon---by exploiting Sentinel-1 C-band backscatter (VH,
VV, RVI) augmented with engineered temporal features (yield lag, rolling
precipitation and temperature means). Evaluated on 11 annual samples
(2015--2025), Random Forest produces independent 2026 forecasts of
\SI{110.78}{\tonne\per\hectare} yield and 1{,}617{,}085 tonnes production.
The real-time operational deployment system (\texttt{predict\_realtime.py})
fetches live weather and SAR data automatically, enabling pre-harvest
forecasting during monsoon without requiring cloud-free optical imagery.

The core contribution of this work is the demonstration that NDVI and SAR
represent complementary rather than competing approaches: NDVI delivers strong
validated accuracy ($R^2 = 0.92$) under clear-sky conditions using large
historical datasets, while SAR provides all-weather operational capability
during the critical monsoon growing season when optical sensors systematically
fail. Together they form a complete, year-round forecasting framework with
substantial practical implications for procurement planning (Food Corporation
of India), weather-index crop insurance (PMFBY), precision input scheduling,
and early drought warning systems.

\subsection{Opportunities for Future Enhancement}

Several promising avenues exist for extending the capabilities of this
framework. Deep learning architectures---particularly Long Short-Term Memory
(LSTM) networks and Temporal Convolutional Networks---can exploit the full
monthly temporal resolution of NDVI and SAR observations, capturing
intra-seasonal phenological dynamics beyond annual aggregations. Geographic
expansion to additional Maharashtra districts and other Indian states would
increase training data volume and provide broader validation of spatial
generalizability. Integration with real-time IoT-based soil moisture and
microclimate sensors would provide ground-truth validation at temporal
resolutions that complement satellite-based observations. Coupling with
process-based crop growth models (WOFOST, DSSAT) via data assimilation
would embed mechanistic agronomic knowledge into the data-driven framework.
Finally, deployment on cloud computing platforms (Google Earth Engine, AWS)
as an automated pipeline would enable operational monthly-updated yield
probability distributions at both district and field scales---realizing the
full potential of the hybrid precision agriculture system developed in this
research.

% -------------------------------------------------------
% ACKNOWLEDGMENTS
% -------------------------------------------------------
\section*{Acknowledgments}

The authors gratefully acknowledge the Department of Computer Science and
Engineering, Indian Institute of Information Technology Dharwad, for
providing computational infrastructure and academic guidance throughout this
research. We thank the NASA Prediction of Worldwide Energy Resources (POWER)
project, the European Space Agency Copernicus program (Sentinel-1 and
Sentinel-2 missions), and Google Earth Engine for providing the open-access
data platforms that enabled this study. Crop production statistics were sourced
from the Ministry of Agriculture and Farmers Welfare, Government of India.

% -------------------------------------------------------
% REFERENCES
% -------------------------------------------------------
\begin{thebibliography}{12}

\bibitem{narimani2026}
M. Narimani et~al., ``Sentinel-2 for Crop Yield Estimation: A Systematic
Review,'' \textit{arXiv preprint arXiv:2603.23779}, Mar.\ 2026.

\bibitem{filgueiras2019}
R. Filgueiras et~al., ``Crop NDVI Monitoring Based on Sentinel~1,''
\textit{Remote Sensing}, vol.~11, no.~12, Art.\ 1441, 2019.

\bibitem{mandal2020}
D. Mandal, V. Kumar, D. Ratha, S. Dey, A. Bhattacharya, J. M.
Lopez-Sanchez, H. McNairn, and Y. S. Rao, ``Dual polarimetric radar
vegetation index for crop growth monitoring using Sentinel-1 SAR data,''
\textit{Remote Sensing of Environment}, vol.~247, Art.\ 111811, 2020.

\bibitem{xu2025}
J. Xu et~al., ``Estimating Sugarcane Biomass via Data Assimilation of
Sentinel-2 LAI into a Crop Model,'' \textit{Int.\ J.\ Appl.\ Earth Obs.\
Geoinf.}, vol.~120, Art.\ 103003, 2025.

\bibitem{chen2023}
J. Chen et~al., ``Crop NDVI Time Series Reconstruction from Sentinel-1 SAR
and Environmental Data,'' \textit{Computers and Electronics in Agriculture},
vol.~207, Art.\ 107685, 2023.

\bibitem{mesas2023}
F. J. Mesas-Carrascosa et~al., ``Enhancing Crop Yield Estimation by
Integrating Sentinel-1 SAR Estimated NDVI,'' \textit{Remote Sensing},
vol.~15, no.~2, Art.\ 296, 2023.

\bibitem{uribe2023}
A. Uribeetxebarria et~al., ``Optimizing Wheat Yield Prediction Integrating
Sentinel-1 and Sentinel-2 with CatBoost,'' \textit{Remote Sensing}, vol.~15,
no.~6, Art.\ 1640, 2023.

\bibitem{aman2024}
K. Amankulova et~al., ``Integrating Sentinel-1, Sentinel-2 and Topographic
Data for Soybean Yield Prediction Using Machine Learning,'' \textit{Advances
in Space Research}, vol.~73, pp.~4052--4066, 2024.

\bibitem{karande2016}
S. V. Karande and P. A. Khadke, ``Climate Variability and Food Security in
Satara District, Maharashtra,'' \textit{European Academic Research}, vol.~4,
no.~4, Jul.\ 2016.

\bibitem{muhammad2026}
A. M. Muhammad et~al., ``A Review of Remote Sensing-Based Crop Yield
Estimation: Machine Learning Techniques and Limitations,'' \textit{Frontiers
in Plant Science}, vol.~17, Art.\ 1742689, 2026.

\bibitem{kameoka2025}
T. Kameoka et~al., ``Rice-Yield Estimation Using Sentinel-2 NDVI and
Interview Data,'' \textit{Scientific African}, vol.~30, Art.\ e03071, 2025.

\bibitem{esa2022}
European Space Agency, \textit{Sentinel-1 SAR GRD: C-band Synthetic Aperture
Radar Ground Range Detected, Log Scaling.} ESA Copernicus Data User Guide,
2022. [Online]. Available:
\url{https://developers.google.com/earth-engine/datasets/catalog/COPERNICUS_S1_GRD}

\bibitem{veloso2017}
A. Veloso et~al., ``Understanding the Temporal Behavior of Crops Using
Sentinel-1 and Sentinel-2-Like Data for Agricultural Applications,''
\textit{Remote Sensing of Environment}, vol.~199, pp.~415--426, 2017.

\end{thebibliography}

\end{document}
